\expandafter\ifx\csname ifdraft\endcsname\relax
   \documentclass[a4j,12pt]{jsreport}
% 数式
\usepackage{amsmath,amsfonts}
\usepackage{bm}

% 画像
\usepackage[dvipdfmx]{graphicx}

% biblatex
\usepackage[style=numeric, sorting=none, maxbibnames=1000]{biblatex}
\addbibresource{graduation_thesis.bib}

%itembox
\usepackage{ascmac}
%breakitemboxのためのpackage
\usepackage{itembkbx}
% \usepackage{tcolorbox} 画像が見えなくなってしまう原因
% \usepackage{framed} 枠

%フォント
\usepackage[ipaex]{pxchfon}

\usepackage[top=35truemm,bottom=30truemm,left=30truemm,right=30truemm]{geometry}
\usepackage{float}
\usepackage{multicol}
\usepackage{multirow}
\usepackage{paralist}
\usepackage{caption}
\usepackage{setspace}
\usepackage{algorithm}
\usepackage{algpseudocode}
\usepackage[dvipdfmx,hidelinks]{hyperref}
\usepackage{pxjahyper}
\usepackage{fancyhdr}
\usepackage{ltablex}
\usepackage{capt-of}
\usepackage{tabularx}
\usepackage{enumitem} % itemizeの行間 これは下になければならない
\usepackage{threeparttable} % 表の脚注


   \setlength{\headheight}{16.5pt}
   \captionsetup{format=hang}
   \begin{document}
   % 上部のヘッダー
\pagestyle{fancy}
\lhead{\rightmark}
\renewcommand{\chaptermark}[1]{\markboth{第\ \thechapter\ 章\ #1}{}}
\fi

\chapter{背景}\label{chap:background}% 物語構造
本章では, 本研究において開発したストーリー創作支援システムが依拠する物語構造について\ref{sec:story_structure}節で述べ, 大規模言語モデルの発展について\ref{sec:llm}節で述べる.

\section{物語構造}\label{sec:story_structure}

\subsection{Transformer}

\expandafter\ifx\csname ifdraft\endcsname\relax
   \fancyhead[LO]{}
   \printbibliography[title=参考文献]
   \end{document}
   \fancyhead[LO]{\rightmark}
\fi